% Section 5. Soft-sensor reframe 2026-11-04 (branch docs/cep-soft-sensor-frame). Estimator design from
% src/softsensor/ (ekf.py, growth.py, evaluate.py) and docs/softsensor/results.md (scripts/softsensor_results.py).
% The validation protocol is moved from csd.tex (sec:validation) of docs/cep-grounded-draft.
% Review round 2. Claim changes: papers/cep/CLAIM_CHANGES.md.
\section{Estimator and baseline}\label{sec:estimator}

This section turns the model of \S\S\,\ref{sec:measurement}--\ref{sec:growth} into the estimator and the references
that RQ3 compares, and fixes the validation protocol before the results.

\subsection{Camera-aided estimator}\label{sec:ekf}

The camera-aided estimator is a discrete extended Kalman filter on $x_k = (\ln L_k, \ln G_k, \theta)^\top$ with the
dynamics~\eqref{eq:growth}. Between measurements it propagates in sub-steps of at most 1~min, so $\Delta t \le 1$~min
in Eq.~\eqref{eq:jacobians} (\code{src/softsensor/ekf.py}). It takes the camera
output $y_k$ (Eq.~\ref{eq:meas}) with the row $C$ and the anchors $z_j$ (Eq.~\ref{eq:anchor}) with the row $C^z$ of
Eq.~\eqref{eq:jacobians}. On a semi-industrial fed-batch crystallizer, extended and unscented Kalman filters estimate
the states accurately, and a disturbance model in the observer recovers performance lost to model error
\citep{Mesbah2011}; the offset $\theta$ plays that role here. The same study compared a moving-horizon estimator, the
alternative if the extended filter proves fragile. The extended filter can fail even without plant--model
mismatch \citep{HaseltineRawlings2005}; with linear measurement rows the only nonlinearity is the size update.

The filter starts at seeding, which the DCS gives. In the evaluation of \secref{sec:validation} the priors come from
the training strike: $\ln L_0$ at its first median and $\ln G_0$ at the log of its mean growth rate, each with
standard deviation 1.0, and $\theta \sim \mathcal{N}(0, 0.16^2)$. The prior standard deviation of $\theta$ is an inference,
$0.16 \approx (0.05^2 + 0.15^2)^{1/2}$: 0.05 for the gain budget (\secref{sec:budget}) and 0.15 as a judgement of the
spread of the selection bias across detectors (\secref{sec:bias}). $Q_G$ is tuned on the training strike over a grid from $10^{-5}$ to
$10$~min$^{-1}$; the optimum is interior, $0.1$~min$^{-1}$ trained on b1 and $0.0316$~min$^{-1}$ trained on b2. The camera updates once per minute while $L \ge 40$~\um, the floor chosen in
\secref{sec:bias}. By Lemma~\ref{lem:scale}, the filter reports two kinds of output. Without an anchor it reports the relative
growth rate $\hat G/\hat L$, the curve shape and the comparison of growth rates between strikes; with an anchor it
also reports the log size $\ln \hat L$, the growth rate $\hat G$ and the growth age $\hat a$ against the nominal
curve. The growth age has two limits: it is ill-conditioned where the nominal curve is flat (on the b1 plateau,
483, 484 and 485.5~\um{} map to 122.5, 137.7 and 160.4~min, so 1\,\% in size moves $\hat a$ by tens of minutes), and
it is undefined above the maximum of the nominal curve (485.5~\um{} for b1, 399.0~\um{} for b2).

For the operator the filter raises a flag on the relative growth rate. With
$d(t) = \ln(\hat G/\hat L)(t) - \ln(\bar G/\bar L)(t)$, the difference to the nominal curve's relative growth rate,
and $\sigma_d$ the posterior standard deviation of $\ln \hat G - \ln \hat L$, a strike is flagged relatively slower
when $d + 1.96\,\sigma_d < -\ln 1.2$ (faster when $d - 1.96\,\sigma_d > \ln 1.2$) on 10 consecutive minutes; the flag
time is the last of these minutes. The offset $\theta$ cancels in $d$, so this flag needs no anchor
(Lemma~\ref{lem:scale}). The same rule applied to $\ln \hat L - \ln \bar L$ gives a size flag, which needs the
anchor. The 20\,\% band, the 1.96 and the 10~min are analysis choices, not tuned.

\subsection{Process-signal baseline and references}\label{sec:baseline}

The baseline answers whether the camera adds information beyond what the process signals already provide. It has
the same growth state without $\theta$ and without the camera, and its growth rate follows the plant's own
supersaturation coefficient: $G = k_g \max(\mathrm{KP} - 1, 0)$, a law linear in $\mathrm{KP} - 1$ (an assumption).
KP is computed exactly as the plant's DCS formula sheet writes it, from Brix, purity and temperature; this form
reproduces the live DCS value KP\_akt of the seed pan, 1.208 against 1.210 (\code{src/softsensor/process.py};
\code{docs/sered/notes_plant_dcs_qa_photos.md}, \code{docs/sered/notes_data_lab.md}). The size starts at
$L_0 = 13.0$~\um, the median of the four chords of the seed-slurry photograph, and $k_g$ is fitted on the training
strike's microscope medians: 16.26~\um/min from b1 and 9.88~\um/min from b2 (\code{scripts/softsensor_process_baseline.py},
\code{docs/softsensor/process_baseline.md}). With the anchor, the predicted size restarts at the anchor. Our group's
supersaturation soft sensor from standard process measurements \citep{Lohani2026} is its predecessor.

No DCS export exists yet for the sampled run of 2026-10-30 (\secref{sec:data}), so the baseline runs in an interim
form: Brix and purity come from the laboratory samples, interpolated linearly, and the temperature is assumed at
70.8~\textdegree C, the median massecuite temperature (TI1006) of the seed pan while boiling in the 2025 DCS data. The
DCS export will replace the laboratory driver by QI1003 at 30~s, the assumed temperature by TI1006 and the operator's
seeding time by the steam cut \pending{DCS export for 2026-10-30, requested from the plant}. It will not remove the
structural limitation of \secref{sec:res-process}: both the laboratory Brix and QI1003 read the whole massecuite. On
the 2026-02-04 export the DCS path of the baseline runs end to end and detects the seeding at 12:31:30, against the
operator's 12:30.

RQ3 also compares the camera-aided estimator with three references without the camera. The \emph{nominal-curve prediction} is the interpolated microscope medians of the
other strike of the same run, read at the same time since seeding; it is what an operator expects from the clock
alone. Two \emph{anchor-only references} pass the nominal curve through the anchor: one shifts it in $\ln L$, the
other shifts it in time, to the growth age at the anchor; before the anchor both equal the nominal curve. With two
strikes, the nominal curve for one is the other, so ``slower than nominal'' for one strike is ``faster'' for the
other by construction.

\subsection{Validation protocol}\label{sec:validation}

The protocol has three levels.

\emph{Detector against hand labels.} Frames are sampled from the crystal-bearing stages of both recordings: the
2026-02-04 PAL recording (15 frames for each of four annotators) and the 2026-10-30 sampled run (12 frames for each of
five annotators). A fifth annotator relabels 15 frames of the first recording for inter-annotator agreement. One frame
per annotator is held out of fine-tuning. Detections are matched to labels one-to-one at intersection-over-union
$\ge 0.5$. Metrics are precision, recall and $F_1$ per detector, the average precision at IoU 0.5 of the fine-tuned
model on the held-out frames, the same metrics between annotators as the attainable ceiling, and the median relative
error and bias of the chord for matched pairs. The edge and watershed detectors are scored on the same frames.

\emph{Camera signal against the microscope.} For the 2026-10-30 run, $y_k$ is compared with $z_j$ at matched times
since seeding. Acceptance (an analysis choice) requires (a) a ratio of camera to microscope median within a
factor of 1.25 at each matched time from 30~min after seeding, which covers the $\pm 10$\,\% gain budget and a margin
for selection; (b) a slope of the log median against time over 30--120~min within $\pm 25$\,\% of the microscope
slope; and (c) a constant camera-to-microscope ratio across the strike: the ratios at all matched times from 30~min
stay within the same factor of 1.25, so that $\theta$ is constant within the strike and no packing bias is present.
Criterion (c) is the key design requirement (\secref{sec:res-robust}). Because the two methods select differently
(\secref{sec:bias}), agreement within the tolerance is necessary, not sufficient.

\emph{Estimator, leave one strike out.} The filter is tuned and given its priors on one strike of the 2026-10-30 run
and tested on the other (\code{src/softsensor/evaluate.py}). Where an anchor is used, it is the first test sample at
$t \ge 30$~min, with an assumed anchor standard deviation of 0.1. Metrics at the held-out microscope times of the test
strike are the log-size RMSE, the growth-rate RMSE against a $\pm 10$~min central difference of the microscope curve,
the RMSE of $G/L$, the growth-age RMSE where it is defined, and the share and median time of flags, against the time
at which the same flag rule fires on the noise-free interpolated microscope curves. The camera-aided estimator passes if it beats the anchor-only references, the nominal curve
and the process-signal baseline on both strikes \pending{acceptance margins, author decision}. Carrying $\theta$ from an anchored strike to
an unanchored one tests whether $\theta$ is shared (Lemma~\ref{lem:scale}(iv)) \pending{b2 with $\theta$ from b1}.
\secref{sec:results} reports what can be run now.
